%% Abstract (configurado para língua inglesa)
\begin{resumo}[Abstract]			% Título do Resumo (Abstract = Resumo em inglês)
\begin{otherlanguage*}{english}		% Língua do texto
The lack of public policies and digital visibility tools has limited access to and recognition of various cultural sites, especially in less explored regions of Brazil. This scenario highlights the need for solutions that promote the recovery and appreciation of local culture, making it more accessible to both residents and visitors. The use of digital technologies, particularly mobile applications, has proven to be a viable and effective strategy to meet this demand. Apps with features such as geolocation, intuitive interfaces, and accessibility can register, disseminate, and make cultural and historical spaces interactive, thus promoting the democratization of cultural access. This study proposes the development of a cultural site registry application, with features focused on inclusion, georeferenced navigation, image presentation, and explanatory texts. The proposal aims to enhance the visibility of the state’s cultural heritage, facilitating engagement and fostering a stronger connection between the population and its history and identity.

\vspace{\onelineskip}
\noindent
\textbf{Key-words}: Digital Culture, Mobile Applications, Geolocation, Cultural Accessibility, Cultural Heritage, Digital Inclusion, Cultural Development, Information Technologies.
\end{otherlanguage*}
\end{resumo}

%% Exemplo de resumo em francês
%\begin{resumo}[Résumé]
% \begin{otherlanguage*}{french}
%    Il s'agit d'un résumé en français.
% 
%   \textbf{Mots-clés}: latex. abntex. publication de textes.
% \end{otherlanguage*}
%\end{resumo}

%% Exemplo de resumo em Espanhol
%\begin{resumo}[Resumen]
% \begin{otherlanguage*}{spanish}
%   Este es el resumen en español.
%  
%   \textbf{Palabras clave}: latex. abntex. publicación de textos.
% \end{otherlanguage*}
%\end{resumo}
% ---